% displayname : Systèmes thermodynamiques ouverts
\documentclass[12pt,fancy]{cours}
\Chapit{2}{1}
\begin{document}
\maketitre{
Ce premier chapitre de thermodynamique étend les principes de la thermodynamique vus pour les systèmes fermés aux systèmes ouverts en régime stationnaire. Les systèmes ouverts sont d'une grande importance pratique et s'étudient à l'aide de diagrammes thermodynamiques : le diagramme des frigoristes $(p,h)$ et entropique $(T,s)$.
}
{
\item Cours Thermodynamique 1 à 4 - 0
\item Documents de cours Thermodynamique 1 - 3, 2 - 3 et 4, 4 - 1 et 2
\item TD Thermodynamique 4
}





%%%%%%%%%%%%%%%%%%%%%%%%%%%%%%%%%%%%%%%%%%%%%%%%%%%%%%%%%%%%%%%%%%%%%%%%%%
\section{Thermodynamique d'un système fermé}

Nous rappelons l'essentiel de la thermodynamique des systèmes fermées vu PCSI.

\subsection{Les principes}


\parag{Systèmes thermodynamiques} 
Considérons un système comme un gaz dans un piston, voir \Fig{piston}. Le gaz est composé d'un nombre macroscopique $N$ d'entités microscopiques : atomes, molécules, ions, etc. En ordre de grandeurs, le nombre de ces entités microscopiques vaut le nombre d’\textsc{Avogadro},
\eqbox{
N \sim \NA = \SI{6.6e23}{}.
}
Il s’agit d’un nombre colossal : en pratique, on ne peut donc pas décrire la position et la vitesse  de chaque entité. On adopte alors une approche \imp{statistique} : on se contente de décrire les propriétés \emph{moyennes} du système. En pratique, on observe que pour beaucoup de systèmes, les propriétés moyennes ne dépendent plus de l’espace ni du temps si l’on attend suffisamment longtemps : on parle d’\imp{équilibre thermodynamique}. Ces grandeurs sont appelées \imp{variables d’état}.
% Les propriétés du gaz peuvent alors êtres décrites par son volume $V$, sa pression $p$ (celle imposée par le piston par équilibre mécanique), sa température $T$ (celle imposée par l'atmosphère), etc.

\begin{definition}[Système thermodynamique fermé]\label{def:}
  Un \imp{système thermodynamique fermé} $\Sigma$ est un corps composé d'un nombre $N$ de particules \imp{fixé} dont les propriétés macroscopiques à l'\imp{équilibre} sont décrites par un petit nombre de \imp{variables d'état}.
  %  : nombre de particules $N$, volume $V$, pression $p$, température $T$, etc.
\end{definition}

Les variables d’états dépendent du système étudié. En plus du nombre de particules $N$ (sans dimension) ou de la quantité de matière $n$ (en mole), qui est fixe pour un système fermé, on peut citer :
\begin{liste}
  \item système hydrostatique (fluide au repos) : pression $p$ (en Pa), température $T$ (en K), volume $V$ (en m$^3$) ;
  \item système hydrodynamique (fluide en mouvement) : vitesse $c$ (en $\SI{}{m.s^{-1}}$), altitude $z$ (en m).
  % \item Système électrique : charge électrique $Q$, tension $u$, intensité $i$, ... ;
  % \item Système mécanique déformable ; force de tension $f$, élongation $L$, ... ;
\end{liste}
\begin{remarque}
  La variable d’état \imp{température $T$} est spécifique à la thermodynamique. Les autres grandeurs sont mécaniques.
\end{remarque}

\begin{figure}[H]
  \centering
  \includegraphics{piston_exemple.pdf}
  \caption{Exemple de système thermodynamique fermé : un gaz de $N$ particules dans un piston de volume $V$. Le gaz est maintenu à une température $T$ et exerce une pression $p$ sur le piston.}
  \label{fig:piston}
\end{figure}


%  On distingue les variables d'états :
% \begin{liste}\item \imp{extensives}, qui sont proportionnelles à la taille du système, i.e. à la masse $m$ ou au nombre de particules $N$, les propriétés du système étant fixées : volume $V$, charge électrique $Q$, élongation $L$.
% \item \imp{intensives}, qui sont indépendantes de la taille du système. On distingue les variables d’états intrinsèquement intensives, et les variables d’états intensives par construction. 
% \begin{itemize}
%   \item La température $T$, la pression $p$, la tension $u$, etc. sont des variables intrinsèquement intensives.
%   \item Le volume massique $v=\frac{V}{m}$, la masse volumique $\rho=\frac{m}{V}$, la densité particulaire $\frac{N}{V}$, etc. sont des variables intensives par construction.
% \end{itemize}

Dans la suite, nous nous intéressons aux \imp{systèmes hydrodynamiques}. Un système thermodynamique est caractérisé par son \imp{équation d’état}, qui est une relation entre l’ensemble de ses variables d’états.
% Un système hydrodynamique possède donc une équation d’état de la forme 
% \begin{equation}
%   f\left( p,T,\frac{N}{V} \right) = 0.
% \end{equation}
Le modèle de fluide le plus simple est celui du \imp{gaz parfait}.

\begin{definition}[Gaz parfait]\label{def:}
  Un \imp{gaz parfait} est un modèle de gaz dont les particules sont \imp{sans interaction}, d’équation d’état
  \begin{equation*}
    % pV = N\kB T = nRT,
    pV = nRT,
  \end{equation*}
   avec  $R=\SI{8.314}{J.K^{-1}.mol^{-1}}$ la constante des gaz parfait.
  %  $\left\{  
  %   \begin{array}{l}
  %     \kB = \SI{1.38e-23}{J.K^{-1}} \text{ la constante de \textsc{Boltzmann}}, \\
  %     R=\kB\NA=\SI{8.314}{J.K^{-1}.mol^{-1}} \text{ la constante des gaz parfait}.
  %   \end{array}
  % \right.$
  \end{definition}
% d'un gaz sont intensives. On peut aussi construire des variables intensives en divisant une variable extensive par la masse $m$ du système, comme le volume massique $v=V/m$. Les grandeurs mécaniques macroscopiques sont aussi des variables intensives, comme la vitesse $c$ et l'altitude $z$ du centre de masse d'un fluide.
% \end{liste}

\parag{Premier principe}

% Étudions le contenu énergétique d’un système hydrodynamique fermé $\Sigma$ évoluant dans le champ de pesanteur. L’énergie se sépare en deux contributions : une contribution macroscopique, et une contribution microscopique.

% \begin{liste}
%   \item \udl{Énergie macroscopique}. 
%   Soit $m_{\Sigma}$ la masse totale du système, dont le centre de masse $G$ se déplace à la vitesse $\vR(G)$ à l’altitude $z(G)$ dans un référentiel $\Rcal$. Ce système évolue dans le champ de pesanteur terrestre $\Vec{g}$. Ce système possède :
%   \begin{itemize}
%     \item une énergie cinétique macroscopique $E_{c, \rm macro} = \frac{1}{2} m_{\Sigma} v_{\! \Rcal}(G)^2$ ;
%     \item une énergie potentielle extérieure $E_{p, \rm ext} = m_{\Sigma} \,g\, z(G)$ ;
%   \end{itemize}
%   La somme de ces deux contributions donne l'\imp{énergie mécanique}, de nature macroscopique,
%   % \begin{equation}
%   \eqbox[]{
%     E_{\rm macro} = E_{c, \rm macro} + E_{p, \rm ext}.
%     }
%   % \end{equation}

%   \item \udl{Énergie microscopique}. Dans le référentiel du centre de masse, les particules $i$ sont de masse $m_i$, vitesse $v_i$, et interagissent entre elles via un potentiel de paire $V(\Vec{r_{i}},\Vec{r_j})$. Le système possède 
%   \begin{itemize}
%     \item une énergie cinétique microscopique $E_{c, \rm micro} = \sum_i \frac{1}{2} m_i v_i^2$ ;
%     \item une énergie potentielle intérieure $E_{p, \rm int} = \sum_{i<j} V(\Vec{r_i},\Vec{r_j})$ ;
%   \end{itemize}
%   La somme de ces deux contributions donne l'\imp{énergie interne}, de nature microscopique,
%   \eqbox[]{
%   U = E_{c, \rm micro} + E_{p, \rm int}.
%   }

% \end{liste}
Pour expliquer les échanges d'énergies observées, on doit admettre qu'en plus de son énergie mécanique macroscopique $E_{\rm m} = E_{\rm c} + E_{\rm p}$, une système thermodynamique fermé possède aussi une énergie d'origine microscopique, appelée énergie interne $U$. Il s'agit d'une \imp{fonction d'état}, i.e. d'une fonction des variables d'états à l'équilibre thermodynamique.
% À l’équilibre thermodynamique, l’énergie interne $U$ est une fonction d’état, i.e. elle ne dépend que des variables d’état du système hydrodynamique $(p,T,N,V)$. De plus, l’énergie interne est extensive.
\begin{definition}[Énergie interne]\label{def:}
  Les propriétés d’un système thermodynamique fermé $\Sigma$ à l’équilibre thermodynamique sont caractérisées par une fonction d’état extensive appelée \imp{énergie interne} $U(T,V)$ (en J).
\end{definition}




% Étudions le contenu énergétique d’un système hydrodynamique fermé $\Sigma$ évoluant dans le champ de pesanteur.
%, que l'on peut caractérisé par son énergie mécan est caractérisé par des grandeurs mécaniques macroscopiques ; en particulier par son \imp{énergie mécanique} $E_{\rm m}=E_{\rm c} + E_{\rm p}$. L'énergie cinétique macroscopique d'un système ayant un mouvement d'ensemble à la vitesse $c$ vaut
% \begin{equation}
% E_{\rm c} = \dfrac{1}{2}mc^2.
% \end{equation}
% L'énergie potentielle macroscopique d'un système en interaction avec l'extérieur, noté $E_{\rm p}$, dépend de la nature de l'interaction. Souvent, il s'agit de la pesanteur, auquel cas $E_{\rm p} = mgz$ si l'altitude du centre de masse est $z$.
L'énergie totale d'un système thermodynamique se décompose donc sous la forme
\eqbox[énergie]{
  E = E_{\rm m} +  U = E_{\rm c} + E_{\rm p} + U.
}
\begin{remarque}
  Dans le cas d’un gaz parfait classique, l’énergie interne est proportionnelle à la température $T$ : 
  \begin{equation}
    U_{\rm GP} = C_V T,
  \end{equation}
  où $C_V = \cte$ est la capacité calorifique à volume constant. On rappelle que 
  \begin{equation}
    \left\{ 
    \begin{array}{l}
      C_V = \frac{3}{2} nR \text{ pour un gaz parfait monoatomique,} \\
      C_V = \frac{5}{2} nR \text{ pour un gaz parfait diatomique.}
    \end{array}
    \right.
  \end{equation}
\end{remarque}



% Un système thermodynamique contient aussi une énergie d'origine microscopique, issue du mouvement et des interactions entre particules à l'échelle atomique, appelée \imp{énergie interne} $U$. Pour un gaz parfait, l'énergie interne ne dépend que de la température : $U=C_V T$, où $C_V$ est une constante. 

% \end{equation}
% Ces trois grandeurs énergétiques sont des \imp{fonctions d'état} : elles dépendent uniquement de l'état du système $(p,T,...)$ à un instant donné, mais pas de l'histoire du système. %Ce sont des fonctions de $V$, $T$, $p$, $z$, etc.
% Ce sont également des grandeurs extensives.

Le premier principe de la thermodynamique stipule la \imp{conservation de l'énergie} : l'énergie d'un système fermé ne peut varier que si de l'énergie est \imp{échangée} avec l'extérieur. Les échanges au cours d'une transformation infinitésimale peuvent prendre deux formes : le \imp{travail} élémentaire $\delta W$ et le \imp{transfert thermique} élémentaire $\delta Q$.
\begin{liste}\item
\udl{Travail} : action mécanique provoquant un mouvement d'ensemble \imp{ordonné} des particules (par exemple lorsqu'un piston pousse un gaz). En particulier, le travail des forces de pression lors d'une transformation élémentaire quasi statique s'écrit
\eqbox{
  \delta W_{\rm p} = -p \,\d V.
  }
% \end{equation}
% où  $\d V$ désigne la variation infinitésimale de volume, et $p_{\rm ext}=p$ si la transformation est quasi-statique (situation d'équilibre mécanique).
\item \udl{Transfert thermique} : échange provenant d'un mouvement \imp{désordonnée} des particules (agitation aléatoire des particules à l'échelle microscopique). Le transfert thermique élémentaire est noté $\delta Q$.

\end{liste}

\begin{theorem}[Premier principe de la thermodynamique -- forme infinitésimale]\label{form:1p}
  La variation infinitésimale d'énergie  d'un système thermodynamique fermé $\Sigma$ vaut
  \begin{equation*}
  \d E = \d E_{\rm c} + \d E_{\rm p} + \d U = \delta W + \delta Q.
  \end{equation*}
  \end{theorem}

  \begin{danger}
    Le travail $W$ et le transfert thermique $Q$ ne dépendent pas que des états initial $A$ et final $B$, mais aussi du chemin suivi dans l'espace des variables d'état par la transformation. C'est pourquoi leurs formes différentielles sont notées avec un $\delta$ et non avec $\d$, comme une différentielle totale.
  \end{danger} 

% \begin{figure}[H]
%   \centering
%   \subfigure[]{\label{fig:}
%   \includegraphics{piston_W.pdf}
%   }\hspace{1cm}
%   \subfigure[]{\label{fig:}
%   \includegraphics{piston_Q.pdf}
%   }
%   \caption{(a) Le travail élémentaire $W$ correspond à un transfert ordonnée d'énergie, sous l'action d'une force~$\F$. La force de pression constante s'exerçant sur un piston de surface $S$ qui se déplace d'une distance $\Delta x$ produit un travail $W = F\Delta x = pS\Delta x = p\Delta V$. (b) Le transfert thermique $Q$ correspond à un transfert désordonné d'énergie.}
%   \label{fig:}
% \end{figure}



% \noindent Considérons une transformation élémentaire entre deux instants voisins $t$ et $t+\d t$, au cours de laquelle un système thermodynamique fermé $\Sigma$ subit une variation infinitésimale d'énergie $\d E = E(t+\d t)- E(t)$. Au cours de cette transformation, le système reçoit de la part de l'extérieur un travail élémentaire $\delta W$ et un transfert thermique élémentaire $\delta Q$.
% \begin{remarque}
% Ces grandeurs ne sont \imp{pas des variations} infinitésimales d'une fonction d'état, et donc ne s'écrivent pas avec le symbole $\d$. Pour indiquer qu'il s'agit de grandeurs élémentaires, on utilise le symbole $\delta$.
% \end{remarque}
%  Appliquons le premier principe à cette transformation élémentaire.


\begin{exemple}{Le cycle Beau de Rochas}
Les moteurs à essence fonctionne sur le cycle Beau de rochas. En raisonnant sur un volume fermé de gaz, le cycle idéal présente les étapes suivantes :%
\begin{liste}
\item $A$-$B$ : compression adiabatique réversible du gaz ;
\item $B$-$C$ : explosion isochore à l’allumage de la bougie (source chaude) ;
\item $C$-$D$ : détente adiabatique réversible  et combustion de l’essence ;
\item $D$-$A$ : ouverture isochore de la soupape, gaz au contact de l’atmosphère (source froide).
\end{liste}
Le gaz est supposé parfait de capacité calorifique $C_V$ constante et on néglige les variations d’énergie mécanique macroscopique. 
\begin{quest}
  \item Exprimer la variation d’énergie $\Delta U$ entre deux états initial $1$ et final $2$.
  \item	
  Déterminer les travaux $W_i$ et le transferts thermiques $Q_i$ à chaque étape $i$ du cycle.
  % \item En déduire le rendement $\eta = -\frac{W_{\rm tot}}{Q_c}$ en fonction des températures $T_A$, $T_B$, $T_C$, et $T_D$.
\end{quest}

\begin{figure}[H]
\centering
\subfigure[]{
\includegraphics[width=0.55\textwidth]{moteur_essence2.jpeg}
}\hspace{0.3cm}
\subfigure[]{
\includegraphics[width=0.39\textwidth]{beau_rochas2.png}
}
\caption{(a) Cycle à quatre temps d’un moteur à essence. (b) Cycle thermodynamique idéal de Beau de Rochas dans le diagramme de Watt $(p,V)$.}
\label{fig:}
\end{figure}
%%%%%%%%%%%%%
\end{exemple}
\examplesol{

 % Déterminons les transferts thermiques pendant les étapes isochores
 \begin{quest}
  \item	
  Astuce : pour une transformation adiabatique, $Q = 0$ et donc $W=\Delta U$; pour une transformation isochore, $W = 0$ et donc $Q=\Delta U$.
  On détermine d'abord le travail $W = -\int p\, \d V$ puis le transfert thermique $Q=\Delta U - W$.
  \begin{liste}
  \item $A$-$B$ : $Q_{AB}=0$ et $W_{AB}=\Delta U = C_V(T_B-T_A)$.
  \item $B$-$C$ : $W_{BC} = 0$ et $Q_{BC} = Q_c = \Delta U = C_V(T_C-T_B)>0$.
\item $C$-$D$ : $Q_{CD}=0$ et $W_{CD}=\Delta U = C_V(T_D-T_C)$.
\item $D$-$A$ : $W_{DA} = 0$ et $Q_{DA} = Q_f = \Delta U = C_V(T_A-T_D)<0$.
\end{liste}
\end{quest}
% Le rendement du moteur vaut
% $$
% \eta = \frac{\text{énergie utile}}{\text{énergie dépensée}} = -\frac{W_{\rm tot}}{Q_c}.
% $$
% Or le premier principe appliqué sur tout le cycle donne $\Delta U = 0 = W_{\rm tot} + Q_c + Q_f$ et donc $-W=Q_c+Q_f$ et $\eta = \frac{Q_c+Q_f}{Q_c} = 1+\frac{Q_f}{Q_c}$, soit
% \eqbox{\eta =  1-\frac{T_D-T_A}{T_C-T_B}.}
}

\begin{experience}{Le moteur Stirling}\label{exp:}
  Nous étudierons le cycle idéal de Stirling en TD. Nous allons illustrer les notions de travail, de transfert thermique, et de rendement sur le moteur de Stirling de la collection. Om rappelle que le cycle idéal comprend quatre étapes : compression isotherme, chauffage isochore, détente isotherme, et refroidissement isochore. Le moteur est alimenté par une source chaude à température $T_c=\SI{270}{\dC}$ par la combustion d'éthanol, la source froide de température $T_f=\SI{60}{\dC}$ étant constitué de l'atmosphère au voisinage du piston. 
  \begin{quest}
    \item	Représenter l'allure du cycle réel dans le diagramme de Watt $(p,V)$. Comparer au cycle idéal.
    \item Déterminer le travail $W$ produit par cycle et la fréquence $f$ du cycle. En déduire la puissance attendue $\mc{P}$ du moteur. Comparer à la puissance électrique mesurée $\mc{P}_{\rm exp}$ en sortie du moteur.
    \item Estimer le transfert thermique $Q$ fourni au gaz par cycle. En déduire le rendement $\eta$ du moteur. Comparer au rendement attendu pour un cycle idéal. On donne l’enthalpie standard de combustion de l’éthanol $\Delta_r H^0 = \SI{-1.24e3}{kJ.mol^{-1}}$.
  \end{quest}
\end{experience}
\examplesol{
\begin{quest}
  \item	ici
  \item On mesure l’aire du cycle et l’on trouve le travail fournit par le moteur $W=\SI{0.14}{J}$. La durée d’un cycle est $\Delta t = \SI{180}{ms}$, ce qui correspond à une fréquence $f=\frac{1}{\Delta t} = \SI{5.6}{Hz}$. La puissance attendue du moteur est alors $\mc{P} = W f = \SI{0.78}{W}$. La puissance électrique mesurée en sortie du moteur est $\mc{P}_{\rm exp} = \SI{0.13}{W}$. On voit qu’il y a des pertes importantes (friction, etc.). 
  \item On peut estimer le rendement du moteur en estimant le transfert thermique $Q$ cédée au gaz par la flamme de la lampe à alcool. On mesure une masse d’éthanol $m=\SI{}{}$
\end{quest}
}

\parag{Second principe}
Le premier principe est un principe de conservation, mais pas un principe d’évolution. En particulier, il ne permet pas de savoir si une transformation $AB$ s’effectue dans la sens $A \rightarrow B$ ou dans le sens $B \rightarrow A$. Pour déterminer le sens d’évolution, nous devons considérer la fonction entropie.

% Un système thermodynamique est aussi caractérisé par une fonction d'état extensive appelée \imp{entropie} et notée $S$. Alors que le premier principe est un principe de conservation, le second principe est un principe d'évolution. Contrairement à l'énergie, l'entropie peut être échangée avec l'extérieur (contribution $S_{\rm e}$) mais aussi créée à l'intérieur du système (contribution $S_{\rm c}$)

\begin{definition}[Entropie]\label{form:1p}
  Les propriétés d’un système thermodynamique fermé $\Sigma$ à l’équilibre thermodynamique sont caractérisées par une fonction d’état extensive appelée \imp{entropie} $S(T,V)$ (en $\SI{}{J.K^{-1}}$). 
\end{definition}

% Considérons une transformation élémentaire entre deux instants voisins $t$ et $t+\d t$, au cours de laquelle un système thermodynamique fermé $\Sigma$ subit une variation infinitésimale d'entropie $\d S = S(t+\d t)- S(t)$. On suppose le système au contact d'un thermostat, de sorte que la surface de $\Sigma$ soit de température homogène $T_{\rm ext}$. Le système subit un transfert thermique élémentaire $\delta Q$. Alors la variation d'entropie se décompose de la façon suivante.

\begin{theorem}[Second principe de la thermodynamique -- forme infinitésimale]\label{form:1p}
  Un système thermodynamique fermé $\Sigma$ est décrit par une fonction d'état extensive $S(T,V)$, appelée \imp{entropie} (en $\SI{}{J.K^{-1}}$). Lors d'une transformation élémentaire, la variation d'entropie vaut
\begin{equation*}
\d S = \delta S_{\rm e} + \delta S_{\rm c} \qquad\text{ avec }\qquad \left\{\begin{array}{ll}
\delta S_{\rm e} = \dfrac{\delta Q}{T_{\rm ext}} & \text{l'entropie échangée}, \\
\delta S_{\rm c} \geq 0 & \text{l'entropie créée}.
\end{array}\right.
\end{equation*}
\end{theorem}
\noindent Le transfert thermique $\delta Q$, et les entropies échangée $\delta S_{\rm e}$
et créée $\delta S_{\rm c}$ dépendent de la transformation.

\begin{definition}
La transformation élémentaire est \imp{réversible} si et seulement si l'entropie créée est nulle, soit
$$
\delta S_{\rm c}=0.
$$
% Une transformation réversible est toujours quasi-statique.
\end{definition}
\noindent Une transformation est réversible lorsque la transformation inverse, qui permet de retourner au même état initial, est possible. Par exemple, la compression infiniment lente d'un gaz dans un piston est réversible ; le contact thermique, même infiniment lent, entre un corps chaud et un corps froid, est irréversible.
\begin{danger}
  Une transformation quasi statique n'est pas forcément réversible ! Par contre, une transformation réversible est toujours quasi statique.
\end{danger}


\begin{exemple}{Le cycle moteur de Carnot}
  Considérons le cycle moteur d'un gaz parfait  constitué de deux isothermes au contact d’une source chaude (température $T_c$) et d’une source froide (température $T_f$), et de deux adiabatiques. 
  
  \begin{quest}
    \item	Déterminer la variation d’entropie $\Delta S$ sur le cycle.  
    \item En déduire que le rendement $\eta$ est minorée par une valeur maximale $\eta_{\rm C}$ correspondant à un cycle réversible, appelé cycle de \imp{Carnot}.
    \item Représenter le cycle de Carnot sur le diagramme de Watt $(p,V)$ en s'aidant de la loi de Laplace.
  \end{quest}
  
  
  \end{exemple}
  \examplesol{
    La variation d’entropie sur le cycle vaut
  $$
  \Delta S = 0 = S_{\rm e} + S_{\rm c} = \frac{Q_c}{T_c} + \frac{Q_f}{T_f} + S_{\rm c} \geq \frac{Q_c}{T_c} + \frac{Q_f}{T_f}.
  $$
  Donc $\frac{Q_f}{Q_c} \leq - \frac{T_f}{T_c}$ et le rendement de ce moteur vaut
  $$
  \eta = 1+\frac{Q_f}{Q_c}\leq 1- \frac{T_f}{T_c} = \eta_{\rm Carnot}.
  $$
  Le rendement maximal $\eta_{\rm Carnot}$ est atteint lorsque toutes les transformations sont réversibles. Le cycle correspondant est appelé \imp{cycle de Carnot}, et représenté en \Fig{carnot}.
  \begin{fig}[carnot]
  \centering
  \includegraphics[width=0.4\textwidth]{carnot.jpeg}
  \capt{\small Cycle de Carnot dans le diagramme de Watt. Les adiabatiques réversibles vérifient, pour un gaz parfait, $p V^\gamma =\cte$ et donc $p \propto 1/V^{\gamma}$. On rappelle que $\gamma>1$ avec $\gamma=\frac{7}{5}=\num{1.4}$ pour un gaz parfait diatomique. Pour les isothermes, $p V =\cte$	 et donc $p \propto 1/V$.}
  \end{fig}
  %%%%%%%%%%%%%

  }
  \medskip
  

Comment relier l'énergie interne $U$ et l'entropie $S$ ? Calculons la variation élémentaire $\d U$ d'un système hydrostatique soumis aux seules forces de pression, en empruntant un chemin fictif réversible. On a $\d S = \delta S_{\rm e}$. De plus, réversible implique $T=T_{\rm ext}$  et $p=p_{\rm ext}$. Le premier principe donne 
\begin{equation}
  \d U = \delta W_{\rm rev} + \delta Q_{\rm rev}  = - p\d V + T \d S.
\end{equation}
Cette relation constitue la relation fondamentale de la thermodynamique. 
\begin{theorem}[Relation fondamentale de la thermodynamique]\label{form:}
  Les variations infinitésimales d'énergie interne $U$ et d'entropie $S$ sont reliées par la relation
  $$
  \d U =  T \d S -p\d V.
  $$
\end{theorem}

\begin{exemple}{Entropie d'un gaz parfait monoatomique}
  Rappeler l'expression de l'énergie interne $U(T)$ d'un gaz parfait monoatomique. En déduire l'expression de l'entropie $S(T,V)$ en s'aidant de la relation fondamentale de la thermodynamique.
\end{exemple}
\examplesol{
  Pour un gaz parfait monoatomique, l'énergie interne est $U = \frac{3}{2} n R T$. En utilisant la relation fondamentale de la thermodynamique, on a
  $$
  \d U = T \d S - p \d V \quad \text{soit} \quad \d S = \frac{\d U}{T} + \frac{p}{T} \d V = \frac{3}{2}n R  \frac{\d T}{T} + nR \frac{\d V}{V}.
  $$
  L'entropie d'un gaz parfait monoatomique est donc
  \eqbox{
    S = \frac{3}{2} nR \ln T + nR \ln V + \cte.
  }
}

\subsection{Classification des machines dithermes}
\begin{definition}[Machine ditherme]\label{def:}
  Une machine \imp{ditherme} opère une transformation \imp{cyclique} sur un système fermé $\Sigma$ au contact d’une \imp{source chaude} et d’une \imp{source froide}.
  % Elle est caractérisée par un \imp{coefficient de performance} $\varepsilon$ défini par
\end{definition}
Au cours de cette transformation cyclique, le système fermé reçoit un travail total $W$. On distingue deux types de machines dithermes, comme représenté sur la \Fig{ditherme}.
\begin{liste}
  \item \imp{Moteur} thermique si $W < 0$ : le système fournit du travail à l’opérateur.
  \item \imp{Récepteur} thermique si $W > 0$ : le système reçoit du travail de l’opérateur.
\end{liste}
Les transferts thermiques avec la source chaude $Q_c$ et la source froide $Q_f$ sont alors contraints par le second principe. On peut représenter l’ensemble des machines dithermes dans un diagramme $(Q_c,Q_f)$, appelé diagramme de Raveau.

\begin{figure}[H]
  \centering
  \subfigure[]{\label{fig:}
  \includegraphics{moteur_ditherme.pdf}
  }\hspace{1cm}
  \subfigure[]{\label{fig:}
  \includegraphics{recepteur_ditherme.pdf}
  }
  \caption{(a) Principe d’un moteur ditherme ($W<0$). Le second principe impose $Q_c >0$ et $Q_f <0$. (b) Principe d’un récepteur ditherme ($W>0$). Le second principe impose $Q_c <0$ et $Q_f >0$.}
  \label{fig:ditherme}
\end{figure}

Chaque type de machines dithermes est caractérisé par un coefficient de performance différent suivant son usage. De façon générale, le coefficient de performance est le rapport 
% \begin{equation}
\eqbox{
  \text{coefficient de performance} = \frac{\text{énergie utile}}{\text{énergie dépensée}}.
  }
% \end{equation}
\begin{table}[h!]
  \centering
  \renewcommand{\arraystretch}{2.5}
  \begin{tabular}{|l|p{0.5\textwidth}|c|}
    \hline
    Type & Exemples & Coefficient de performance \\
    \hline
    \imp{Moteur thermique} & moteur Diesel, Otto-Beau de Rochas, Stirling & $\displaystyle \eta=-\frac{W}{Q_c}$ \\[3mm]
    \hline
    \imp{Machine frigorifique} & réfrigérateur, congélateur, climatiseur & $\displaystyle e=\frac{Q_f}{W}$ \\[3mm]
    \hline
    \imp{Pompe à chaleur} & pompe à chaleur, climatiseur & $\displaystyle \mathrm{COP}=\frac{Q_c}{W}$ \\[3mm]
    \hline
  \end{tabular}
  \caption{Classification des machines dithermes et leurs coefficients de performance.}
\end{table}

\begin{prop}[Théorème de Carnot]\label{prop:}
  Le coefficient de performance d'une machine ditherme réelle est toujours \imp{inférieur} à celui de la machine \imp{réversible} équivalente, appelé coefficient de \imp{Carnot}.
\end{prop}

\begin{exemple}{Coefficients de Carnot}
  \begin{quest}
    \item	Déterminer l’efficacité frigorifique maximale $e_{\rm C}$ d’un réfrigérateur fonctionnant entre une source froide à température $T_f=\SI{4}{\dC}$ et une source chaude à température $T_c=\SI{24}{\dC}$.
    \item Déterminer le coefficient de performance maximal ${\rm COP}_{\rm C}$ d’une pompe à chaleur fonctionnant entre une source froide à température $T_f=\SI{4}{\dC}$ et une source chaude à température $T_c=\SI{24}{\dC}$.
  \end{quest}
\end{exemple}
\examplesol{}

%%%%%%%%%%%%%%%%%%%%%%%%%%%%%%%%%%%%%%%%%%%%%%%%%%%%%%%%%%%%%%%%%%%%%%%%%%
\section{Thermodynamique d'un système ouvert}

\subsection{Système thermodynamique ouvert}

\parag{Définition}
Dans la plupart des machines industrielles, un fluide circule entre différents organes de la machine. Le volume de fluide que contient un organe donné constitue un système ouvert, puisque du fluide entre constamment en amont et ressort en aval.

\begin{definition}\label{def:systeme_ouvert}
Un \imp{système thermodynamique ouvert} $\So$ est un volume délimité par une surface de contrôle fixe traversée par de la matière.
\end{definition}
\noindent Conformément au programme, nous étudions uniquement des écoulements unidimensionnels, pour lesquels les grandeurs physiques sont uniformes sur les sections d'entrée et de sortie.

%%%%%%%%%%%%%%%%%%%%%%%%%%%%%
\begin{figure}[h!]
\centering
\includegraphics[scale=1]{syst_ouvert2.pdf}
\caption{Un système ouvert $\So$ (en vert) est délimité par une surface de contrôle fixe traversée par de la matière (les masses entrante $\dM_{\rm 1}$ et sortante  $\dM_{\rm 2}$). En amont de $\So$ entre du fluide à la pression~$p_{\rm 1}$, vitesse $c_{\rm 1}$, volume massique $v_1$, etc. En aval de $\So$ sort du fluide à la pression $p_{\rm 2}$, vitesse $c_{\rm 2}$, volume massique $v_2$, etc.}
\label{fig:ouvert}
\end{figure}
%%%%%%%%%%%%%

%%%%%%%%%%%%%%%%%%%%%%%%%%%%%%
%\begin{figure}[b!]
%\centering
%\includegraphics[scale=1]{echangeur.pdf}
%\caption{Un système ouvert $\So$ (en tireté noir) est délimité par une surface de contrôle fixe traversée par de la matière (les masses entrante $\d m_{\rm e}$ et sortante  $\d m_{\rm s}$). En amont de $\So$ entre du fluide à la pression~$p_{\rm e}$, vitesse $c_{\rm e}$, etc. En aval de $\So$ sort du fluide à la pression $p_{\rm s}$, vitesse $c_{\rm s}$, etc.}
%\label{fig:ouvert}
%\end{figure}
%%%%%%%%%%%%%%



Considérons l'écoulement du fluide à travers la surface de contrôle de $\So$ entre deux instants voisins $t$ et $t + \d t$. Soit  $\dM_{\rm 1}$ la masse entrante et $\dM_{\rm 2}$ la masse sortante. Le bilan de masse appliqué au système ouvert $\So$ s’écrit simplement (on compte ce qui entre et ce qui sort) \begin{equation}
m_{\Sigma_{\rm{o}}}(t+\d t) = m_{\Sigma_{\rm{o}}}(t)+\dM_{\rm 1} - \dM_{\rm 2}.
\end{equation}
Bien souvent, le flux de matière est constant à l'entrée et à la sortie : les masses entrante $\dM_{\rm 1} = {\Dm}_1\d t$ et sortante $\dM_{\rm 2} = {\Dm}_2 \d t$ sont proportionnelles à la durée $\d t$ du bilan. Les constantes ${\Dm}_1$ et ${\Dm}_2$ se mesurent en $\SI{}{kg.s^{-1}}$ et s'appellent débits massiques (respectivement entrant et sortant).

\begin{definition}
On appelle \imp{débit massique} la masse $\delta m$ traversant une section droite par unité de temps $\d t$, soit
$$
\Dm = \frac{\delta m}{\d t}
$$
\end{definition}

\parag{Régime stationnaire} Supposons de plus que l'écoulement est stationnaire (permanent). L'état du fluide contenu dans $\So$, et en particulier sa masse, ne dépend alors pas du temps, soit
\begin{equation}
m_{\Sigma_{\rm{o}}}(t+\d t) = m_{\Sigma_{\rm{o}}}(t) ;
\end{equation}
et ce bien que le fluide  soit constamment renouvelé à l'intérieur de la surface de contrôle de $\So$. Donc les masses entrante et sortante et les débits massiques correspondant, sont égaux.


\begin{prop}\label{prop:}
Dans un écoulement \imp{stationnaire} de débit massique $\Dm$ étudié pendant une durée $\d t$,  on a
\begin{equation*}
\dM_{\rm 1} = \dM_{\rm 2} = \d m = \Dm \d t.
\end{equation*}
\end{prop}
\noindent Il y a donc conservation du débit massique en régime stationnaire.



\parag{Système fermé associé} Pour appliquer le premier principe de la thermodynamique, et plus généralement effectuer le bilan d'une grandeur extensive dans un système ouvert en régime stationnaire, il faut se ramener à un système fermé. Pour un système ouvert $\So$ donné étudiée entre deux instants $t$ et $t+\d t$, nous définissons le système thermodynamique fermé $\Sf$ tel que :
\begin{liste}\item
\udl{Instant $t$} : $\Sf(t)$ se compose de $\So$ et de la masse $\delta m_{\rm 1}$ entrant en amont de $\So$ entre les instants $t$ et $t+\d t$ (en tireté rouge sur la \Fig{ferme});
\item
\udl{Instant $t+\d t$} : $\Sf(t+\d t)$ se compose de $\So$ et de la masse $\delta m_{\rm 2}$ sortant en aval de $\So$ entre les instants $t$ et $t+\d t$ (en tireté bleu sur la \Fig{ferme}).
%%%%%%%%%%%%%%%%%%%%%%%%%%%%%
\begin{figure}[h!]
\centering
\includegraphics[scale=1]{syst_ferme.pdf}
\caption{Le système fermé $\Sf$ associé au système ouvert $\So$ est délimité par $\So$ + masse entrante $\dM_{\rm 1}$ à l'instant $t$, et de $\So$ + masse sortante  $\dM_{\rm 2}$ à l'instant $t+\d t$.
% En amont de $\So$ entre du fluide à la pression~$p_{\rm 1}$, vitesse $c_{\rm 1}$, etc. En aval de $\So$ sort du fluide à la pression $p_{\rm 2}$, vitesse $c_{\rm 2}$, etc.
}
\label{fig:ferme}
\end{figure}
%%%%%%%%%%%%%

%Soit $E=U+E_{\rm{c}}+E_{\rm{p}}$ l'énergie totale (somme de l'énergie interne microscopique et de l'énergie macroscopique c'est-à-dire l'énergie mécanique), alors
%\begin{equation}
%\begin{cases}
%E_{\Sigma_{\rm{f}}}(t)=E_{\Sigma_{\rm{o}}}(t)+\delta E_{\rm{e}},\\
%E_{\Sigma_{\rm{f}}}(t+\d t)=E_{\Sigma_{\rm{o}}}(t+\d t)+\delta E_{\rm{s}}.
%\end{cases}
%\end{equation}

\end{liste}
Le bilan de masse effectué cette fois-ci sur le système fermé $\Sf$ s’écrit
\begin{equation}
\begin{cases}
m_{\Sigma_{\rm{f}}}(t)=m_{\Sigma_{\rm{o}}}(t)+\dM_{\rm{1}},\\
m_{\Sigma_{\rm{f}}}(t+\d t)=m_{\Sigma_{\rm{o}}}(t+\d t)+\dM_{\rm{2}}.
\end{cases}
\end{equation}
Par définition d’un \imp{système fermé}, nous avons
% \begin{equation}
\eqbox[]{
  m_{\Sigma_{\rm{f}}}(t) = m_{\Sigma_{\rm{f}}}(t+\d t),
  }
% \end{equation}
et donc retrouve bien que $m_{\Sigma_{\rm{o}}}(t+\d t)= m_{\Sigma_{\rm{o}}}(t)+\dM_{\rm 1} - \dM_{\rm 2}$. En régime \imp{stationnaire}, on a de plus $m_{\Sigma_{\rm{o}}}(t) = m_{\Sigma_{\rm{o}}}(t+\d t)$ car les propriétés du volume de contrôle $\So$ ne varient pas avec le temps. 
% Donc les masses entrante et sortante sont égales, et les débits massiques correspondants sont égaux.


\begin{methode}[Effectuer un bilan sur un système thermodynamique ouvert]
  \item	Définir un volume de contrôle $\So$ fixe traversé par de la matière. Définir un \imp{système fermé} $\Sf$ associé à $\So$ aux instants $t$ et $t+\d t$ en incluant les masses entrante $\delta m_1$ et sortante $\delta m_2$.
  \item Exprimer la \imp{grandeur extensive} $X_{\Sf}$ aux instants $t$ et $t+\d t$ en fonction des masses $\delta m_1$ et $\delta m_2$.
  \item Utiliser la \imp{stationnarité} de l'écoulement pour établir que $X_{\So}$ est constant et $\delta m_1 = \delta m_2 = \d m$.
  \item Utiliser le caractère \imp{conservatif} de $X_{\rm \Sf}$ ou toute autre loi pour exprimer la variation de $X_{\rm \Sf}$.
  \item Diviser par $\d m$ (ou $\d t$)  pour faire apparaître la grandeur \imp{massique} $x$ (ou le \imp{débit} $\Dm$). 
\end{methode}

\subsection{Premier principe industriel}


\parag{Variation d’énergie}
Effectuons de même un bilan d’énergie totale $E=U+E_{\rm{c}}+E_{\rm{p}}$ (somme de l'énergie interne microscopique et de l'énergie mécanique macroscopique) sur le système fermé~$\Sf$. 
\begin{enumerate}
  \item Nous avons déjà défini le volume de contrôle $\So$ et le système fermé $\Sf$ associé à $\So$.
  \item L'énergie du système fermé vaut
  \begin{equation}
    \begin{cases}
      E_{\Sigma_{\rm{f}}}(t)=E_{\Sigma_{\rm{o}}}(t)+\delta E_{\rm{1}},\\
      E_{\Sigma_{\rm{f}}}(t+\d t)=E_{\Sigma_{\rm{o}}}(t+\d t)+\delta E_{\rm{2}},
    \end{cases}
  \end{equation}
où $\delta E_{\rm{1}}$ et $\delta E_{\rm{2}}$ désignent les énergies entrante et sortante, respectivement. Le système n’est pas isolé, de sorte que l’énergie n’est pas constante. On note $\d E = E_{\Sigma_{\rm{f}}}(t+\d t) - E_{\Sigma_{\rm{f}}}(t)$ la variation d’énergie. 
\item 
Lorsque le régime \imp{stationnaire} est atteint, on a de plus $E_{\Sigma_{\rm{o}}}(t+\d t) = E_{\Sigma_{\rm{o}}}(t)$, et donc
\begin{equation}
\d E = \delta E_{\rm 2} - \delta E_{\rm 1} = \JRouge{\underbrace{(\delta U_{\rm 2} - \delta U_{\rm 1})}_{\d U}} + \JVert{\underbrace{(\delta E_{\rm c, 2} - \delta E_{\rm c, 1})}_{\d E_c}} + \JBleu{\underbrace{(\delta E_{\rm p, 2} - \delta E_{\rm p, 1})}_{\d E_p}}:
\end{equation}
 la variation d’énergie de $\Sf$ provient de l’énergie du fluide échangé avec l’extérieur de $\So$.  
%  Elle se décompose en trois parties : $\JRouge{\d U = \delta U_{\rm 2} - \delta U_{\rm 1}}$, $\JVert{\d E_{\rm c} = \delta E_{\rm c, 2} - \delta E_{\rm c, 1}}$, et $\JBleu{\d E_{\rm p} = \delta E_{\rm p, 2} - \delta E_{\rm p, 1}}$.
 \item Nous pouvons enfin appliquer le \imp{premier principe} au système fermé $\Sf$ (on ne peut appliquer les principes de la thermodynamique qu’à un système fermé!). Il vient
 \begin{align}\label{eq:Ee_ouvert}
 \d E = \d U + \d E_c + \d E_p & = \delta W + \delta Q = \delta W_p + \delta W_u + \delta Q,
 \end{align}
 où nous décomposons le travail $\delta W$ en travail des forces de pression $\delta W_p = \delta W_{p,1} + \delta W_{p,2}$, somme du travail moteur exercé par la pression en amont ($1$) et du travail résistant exercé par la pression en aval ($2$), et en travail utile $\delta W_u$, travail fourni au fluide par les pièces mobiles.
 \item 
Nous allons exprimer chacun de ces termes énergétiques en fonction des énergies massiques.
\begin{table}[h!]
  \centering
  \renewcommand{\arraystretch}{1.2}
  \begin{tabular}{|c|c| c| c| c| c| c| c| c|}
    \hline
    Énergie (J) & $U$ & $E_c=\frac{1}{2}mc^2$ & $E_p=mgz$ & $W_p=-p\d V$ & $W_{p,1}>0$ & $W_{p,2}<0$ & $W_u$ & $Q$ \\
    \hline
    ... massique (\SI{}{J.kg^{-1}}) & $u$ & $e_c = \frac{1}{2}c^2$ & $e_p=gz$ & $w_p=-pv$ & $w_{p,1}=p_1v_1$ &  $w_{p,1}=-p_2v_2$ & $w_u$ & $q$ \\
    \hline
  \end{tabular}
  % \caption{}
  % \label{}
\end{table}

%  \begin{liste}
%   \item La variation d'énergie interne $\JRouge{\d U = \delta U_{\rm 2} - \delta U_{\rm 1} = \d m(u_2-u_1)}$,   où $u$ désigne l'énergie interne massique ;
%   \item La variation d'énergie cinétique $\JVert{\d E_c = \delta E_{\rm c, 2} - \delta E_{\rm c, 1} = \d m(e_{c,2}-e_{c,1})}$, où
%   %\begin{equation}
%   $ e_c=\frac{1}{2}c^2$
%   % \end{equation}
%   désigne l'énergie cinétique massique ;
%   \item La variation d'énergie potentielle $\JBleu{\d E_p = \delta E_{\rm p, 2} - \delta U_{\rm p, 1} = \d m(e_{p,2}-u_{p,1})}$,   où $e_p$
%   désigne l'énergie potentielle massique (de pesanteur en général, soit $e_p=gz$).
On obtient
  \begin{align}\label{eq:Ee_ouvert}
    \d m (u_2-u_1) +  \d m (e_{c,2}-e_{c,1}) + \d m(e_{p,2}-u_{p,1})  = \d m(p_1v_1-p_2v_2) + w_{\rm u} \d m + q \d m.
    \end{align}

On divise par la masse $\d m$ et on regroupe les termes en entrée et les termes en sortie,
\begin{equation}
(u_2 + p_2v_2 + e_{c,2} + e_{p,2}) - (u_1 + p_1 v_1 +e_{c,1} + u_{p,1})  = \left[u + pv + e_{c} + e_{p} \right]_1^2 =  w_{\rm u} + q.
\end{equation}
On voit apparaître l'\imp{enthalpie massique} $h = u+ pv$.
%  En exprimant $e_{c}  = \frac{1}{2}c^2$ et en assimilant $e_{\rm p}$ à l’energie potentielle de pesanteur (seul cas que vous rencontrerez), on trouve finalement,
\end{enumerate}

\begin{theorem}[Premier principe industriel]\label{form:1principe}
La variation d’\imp{énergie massique} (en \SI{}{J.kg^{-1}}) entre l’entrée $1$ et la sortie $2$  d'un système thermodynamique \imp{ouvert} s’écrit en régime \imp{stationnaire}  :
\begin{equation*}
\left[ h + g z + \frac{1}{2}c^2\right]_1^2  = w_u + q.
\end{equation*}
\end{theorem}

% \begin{remarque}
Le premier principe peut aussi s'exprimer en terme de puissances (en $\SI{}{W} = \SI{}{J.s^{-1}}$) en divisant \Eq{Ee_ouvert} par la durée $\d t$. En posant $\Phi_{\rm{th}}=\frac{\delta Q}{\d t}=\Dm q$ le flux de chaleur et $P_{\text{u}}=\frac{\delta W_{\rm u}}{\d t} = \Dm w_{\rm u}$ la puissance utile, 
% \begin{equation}
% \begin{split}
\eqbox{
  \Dm\,\left[h +g z+\frac{1}{2}c^{2}\right]=P_{\rm{u}}+\Phi_{\rm{th}}.
  }
% \end{split}
% \end{equation}
%Attention à l'homogénéité! Dans l'expression ci-dessus, chaque terme est homogène à des $\SI{}{J.s^{-1}}=\SI{}{W}$.
% \end{remarque}

\subsection{Second principe industriel}

On raisonne de même que précédemment pour un bilan l’entropie.
%  On applique le second principe de la thermodynamique au système fermé $\Sf$ entre $t$ et $t+\d t$. 
\begin{enumerate}
  \item Nous avons déjà défini le volume de contrôle $\So$ et le système fermé $\Sf$ associé à $\So$.
  \item L'entropie du système fermé vaut
  \begin{equation}
    \begin{cases}
      S_{\Sf}(t) =S_{\Sigma_{\rm{o}}}(t)+\delta S_{\rm{1}},\\
      S_{\Sf}(t+\d t ) =S_{\Sigma_{\rm{o}}}(t+\d t)+\delta S_{\rm{2}},
    \end{cases}
  \end{equation}
où $\delta S_{\rm{1}}$ et $\delta S_{\rm{2}}$ désignent les énergies entrante et sortante, respectivement. On note $\d S = S_{\Sigma_{\rm{f}}}(t+\d t) - S_{\Sigma_{\rm{f}}}(t)$ la variation d'entropie. 
\item 
Lorsque le régime \imp{stationnaire} est atteint, on a de plus $S_{\Sigma_{\rm{o}}}(t+\d t) = S_{\Sigma_{\rm{o}}}(t)$, et donc
\begin{equation}
\d S = \delta S_{\rm 2} - \delta S_{\rm 1}.
\end{equation}
 la variation l'entropie de $\Sf$ provient de l'entropie du fluide échangé avec l’extérieur de $\So$.  
 \item Nous pouvons enfin appliquer le \imp{second principe} au système fermé $\Sf$. Il vient
 \begin{align}\label{eq:Ee_ouvert}
 \d S = & = \delta S_{\rm e} + \delta S_{\rm c},
 \end{align}
 où $\delta S_{\rm e} = \frac{\delta Q}{T_{\rm ext}}$ est l'entropie échangée, et $\delta S_{\rm c}\geq 0$ l'entropie créée.
 \item 
Nous allons exprimer chacun de ces termes entropiques en fonction des entropies massiques.
\begin{table}[h!]
  \centering
  \renewcommand{\arraystretch}{1.5}
  \begin{tabular}{|c|c| c| c| c| c| c| }
    \hline
    Entropie (\SI{}{J.K^{-1}}) & $S$ &  $\delta S_{\rm e} = \frac{\delta Q}{T_{\rm ext}}$ & $\Delta S_{\rm c}$ \\
    \hline
    ... massique (\SI{}{J.K^{-1}.kg^{-1}}) & $s$   & $s_{\rm e} = \frac{q}{T_{\rm ext}}$ & $s_{\rm c} $ \\
    \hline
  \end{tabular}
  % \caption{}
  % \label{}
\end{table}

On obtient
  \begin{align}\label{eq:Se_ouvert}
    \d m (s_2 - s_1) = \d m s_{\rm e}  + s_{\rm c}.
    \end{align}
On divise par la masse $\d m$ et on trouve le second principe industriel.
\end{enumerate}


\begin{theorem}[Second principe industriel]\label{form:2principe}
La variation d’\imp{entropie massique} (en $\SI{}{J.K^{-1}.kg^{-1}}$) entre l’entrée $1$ et la sortie $2$ d'un système thermodynamique \imp{ouvert} s’écrit en régime \imp{stationnaire}  :
\begin{equation*}
[ s]_1^2=\frac{q}{T_{\rm{ext}}}+s_{\text{c}},
\end{equation*}
où $s_{\rm c} = 0$ si la transformation est réversible, et  $s_{\rm c} > 0$ sinon.

\end{theorem}
% \begin{remarque}
Le taux de variation d’entropie (en  $\SI{}{J.K^{-1}.s^{-1}}$) s'obtient en divisant \Eq{Se_ouvert} par $\d t$, ce qui donne
% \begin{equation}
\eqbox{
  \Dm \,[ s]_1^2=\frac{\Phi_{\rm{th}}}{T_{\rm{ext}}}+\frac{\delta s_{\text{c}}}{\d t}.
  }
% \end{equation}
% \end{remarque}



\subsection{Quelques organes de machine industrielle}

Dans toutes la suite, nous considérons des composants de machine industrielle dans lesquels la variation d'énrgie cinétique et potentielle est négligeable. Nous aurons donc 
\begin{equation}
  \left[ \frac{1}{2}c^2 + gz \right]_1^2 = 0.
\end{equation}
Cela exclut les tuyères (ou $\frac{1}{2}c^2$ varie significativement) et les barrages (ou $gz$ varie significativement).

% %%%%%%%%%%%%%%%%%%%%%%%%%%%%
% \begin{figure}[H]
% \centering
% % \begin{minipage}{0.23\textwidth}
% \subfigure[]{\label{fig:compresseur}
% \includegraphics[scale=1]{compresseur.pdf}
% }
% \hfill
% \subfigure[]{\label{fig:turbine}
% \includegraphics[scale=1]{turbine.pdf}
% }
% % \end{minipage}
% \hfill
% % \begin{minipage}{0.22\textwidth}
% \centering
%   \subfigure[]{\label{fig:detendeur}
%   \includegraphics[scale=1]{laminage.pdf}
%   }
%   \hfill
%   \subfigure[]{\label{fig:echangeur}
%   \includegraphics[scale=1]{evaporateur.pdf}
%   }
%   \subfigure[]{\label{fig:}
%   \includegraphics[scale=1]{condenseur.pdf}
%   }
% \caption{Les principaux organes des machines industrielles, où sont indiqués la variation de pression et le transfert d'énergie. (a) Compresseur. (b) Turbine. (c) Détendeur calorifugé. (d) Échangeur thermique.
% }
% \label{fig:organes}
% \end{figure}
% %%%%%%%%%%%%%




\parag{Compresseur, pompe et turbine}
Le but d’un compresseur (\Fig{compresseur}), d'une pompe ou d’une turbine (\Fig{turbine}) est de faire varier la pression en fournissant ou soutirant un travail utile. Un compresseur et une pompe augmentent la pression du fluide ($p_2>p_1$) tandis qu’une turbine diminue la pression du fluide ($p_2<p_1$). Un compresseur agit sur un gaz, tandis qu'une pompe agit sur un liquide.

La transformation du fluide dans un compresseur, une pompe ou une turbine est \imp{adiabatique réversible}. Adiabatique signifie $q=0$. Réversible signifie $s_c=0$. D’après le second principe industriel, $s_{2} - s_{\rm 1} = \frac{q}{T_{\rm ext}} + s_{\rm c} = 0$. La transformation est donc isentropique. Le premier principe industriel donne alors $h_2-h_1 = w_{\rm u}.$
On a $w_{\rm u}>0$ pour un compresseur (l'opérateur exerce un travail pour comprimer le fluide) et $w_{\rm u}<0$ pour une turbine (l'opérateur récupère du travail des pales mises en rotation par le fluide).

\begin{prop}[Compresseur et turbine]\label{prop:}
La transformation du fluide dans un compresseur ou une turbine est \imp{adiabatique réversible}, de travail utile massique reçu $w_u = \Delta h$ ($>0$ pour un compresseur ou une pompe, et $<0$ pour une turbine).
% \begin{fig}[label]
\begin{figure}[H]
  \centering
  % \includegraphics{}
  \subfigure[Compresseur]{\label{fig:}
  \includegraphics[scale=1]{compresseur.pdf}
  }\hspace{0.25\textwidth}
  \subfigure[Turbine]{\label{fig:}
  \includegraphics[scale=1]{turbine.pdf}
  }
  % \caption{}
  % \label{fig:}
\end{figure}
\end{prop}


\parag{Détendeur} Un détendeur calorifugé, ou vanne de détente calorifugée (\Fig{detendeur}), opère un laminage, i.e. une chute de pression ($p_2<p_1$) sans travail utile ($w_{\rm u}=0$) ni transfert thermique ($q=0$). Or le premier principe industriel donne $h_2-h_1 = w_{\rm u} + q = 0$.
Par contre, la transformation est irréversible dans un cycle réel ($s_{\rm c}>0$). En effet, un fluide réel ne peut pas évoluer à la fois à $h$ et $s$ constants.
\begin{prop}[Détendeur]\label{prop:}
La transformation du fluide dans un détendeur calorifugé est \imp{isenthalpique}.
\begin{figure}[H]
  \centering
  % \subfigure[]{\label{fig:detendeur}
  \includegraphics[scale=1]{laminage.pdf}
  % }
  % \caption{}
  % \label{fig:}
\end{figure}
\end{prop}

\parag{Échangeur thermique} Le but d’un échangeur thermique est de chauffer et de refroidir le fluide sans travail utile ($w_{\rm u}=0$). Un échangeur est donc au contact d'une source, chaude ou froide. Le premier principe donne alors $h_2-h_1=q.$ On distingue des échangeurs particuliers.
\begin{liste}
  \item
  Un \imp{évaporateur} est un échangeur thermique particulier dans lequel un mélange liquide-vapeur est entièrement vaporisé.
  %  Le fluide est alors au contact d'une \imp{source chaude}.
  \item Un \imp{condenseur} est un échangeur thermique particulier dans lequel un mélange liquide-vapeur est entièrement liquéfié. Par extension, on appelle aussi condenseur tout échangeur thermique dans lequel le fluide est refroidi. 
  % Le fluide est alors au contact d'une \imp{source froide}.
  \item Un générateur de vapeur est un échangeur thermique qui chauffe le liquide, le vaporise, puis chauffe encore la vapeur ainsi formée.
\end{liste}
 \begin{prop}[Échangeur thermique]\label{prop:}
 La transformation du fluide dans un évaporateur ou un condenseur est \imp{isobare}, de transfert thermique massique reçu $q = \Delta h$ ($>0$ pour un évaporateur et $<0$ pour un condenseur).
 \begin{figure}[H]
  \centering
  \setcounter{subfigure}{0}
  \subfigure[Évaporateur]{\label{fig:echangeur}
  \includegraphics[scale=1]{evaporateur.pdf}
  }\hspace{0.25\textwidth}
  \subfigure[Condenseur]{\label{fig:}
  \includegraphics[scale=1]{condenseur.pdf}
  }
 \end{figure}
 \end{prop}
\begin{danger}
  Attention, le signe de $q$ et donc la nature de l'échangeur thermique ne préjuge pas de la nature de la source de chaleur !
  Pour un cycle \imp{réfrigérant} ou une \imp{pompe à chaleur} (PAC), l'évaporateur est au contact de la source froide, et le condenseur est au contact de la source chaude.
  Pour un cycle \imp{moteur}, l'évaporateur est au contact de la source chaude, et le condenseur est au contact de la source froide.
\end{danger}
\noindent
Nous avons représenté un exemple de cycle industriel simple : le cycle frigorifique des réfrigérateur, \Fig{frigo}.


  \begin{figure}[H]
  \centering
  \subfigure[]{\label{fig:cycle_fonct}
  \includegraphics[width=0.42\textwidth]{cycle_frigo.pdf}
  }
  \subfigure[]{\label{fig:frigo}
  \includegraphics[width=0.45\textwidth]{frigo_bis}
  }
  \caption{(a) Exemple de cycle de fonctionnement industriel : le cycle frigorifique. (b) fonctionnement d’un cycle frigorifique dans un réfrigérateur : le fluide frigorifique est comprimé et échauffé à l’état gazeux puis refroidi au contact de l’air derrière le réfrigérateur dans un condenseur, puis partiellement liquéfié et refroidi dans un détendeur ; enfin le fluide s’échauffe au contact des compartiments du congélateur et du réfrigérateur.~[Source : \emph{La physique par les objets du quotidien}].
  % avec thermostat.
  %(e) Échangeur thermique avec autre fluide.
  }
  \label{fig:frigo}
  \end{figure}
  %%%%%%%%%%%%%



\section{Diagrammes thermodynamiques}


Un diagramme thermodynamique présente l’état d’un fluide donné par deux paramètres, l'un étant intensif, l'autre étant une fonction d'état massique : ($p,h$) pour le diagramme des frigoristes (\Sec{ph}) et $(T,s$) pour le diagramme entropique (\Sec{Ts}).
Les diagrammes thermodynamiques sont issus de données expérimentales, et permettent donc d'aller plus loin que des modèles simplistes (gaz parfait, liquide incompressible et indilatable). Sur le diagramme on représente aussi un faisceau de courbes renseignant sur les autres variables d'état du système.

\begin{danger}
  N'utilisez pas la loi du gaz parfait lorsqu'on vous donne un diagramme thermodynamique réel ! Le but est justement d'utiliser des données expérimentales et non d'approximer le comportement du fluide par une loi simple.
\end{danger}


\subsection{Diagramme $(p,h)$}
\label{sec:ph}

Un diagramme des frigoristes représente la pression $p$ en ordonnée en échelle \imp{logarithmique} et l’enthalpie $h$ en abscisse en échelle linéaire. Un exemple est le  diagramme de l’eau (réfrigérant R718) de la \Fig{frigoristes_eau}. Lorsque vous voyez un diagramme $(p,h)$, la première chose à faire est d'identifier correctement les courbes représentées.

%Il existe un diagramme tabulé pour chaque fluide frigorigène. Dans la suite nous étudierons le R718, plus connu sous le nom ... d'eau.

\parag{Courbe de saturation} Le diagramme contient une courbe en cloche : la courbe de saturation. Sous la courbe de saturation, le fluide est diphasé : on a coexistence liquide-vapeur $L+V$. De part et d'autre de la courbe de saturation, le fluide est sous une seule phase : liquide $L$ aux basses enthalpies, et vapeur $V$ aux hautes enthalpies. Au dessus de la courbe de saturation, dont le sommet est le point critique $C$, le fluide est en phase supercritique. La partie gauche de la courbe de saturation est la courbe d’ébullition (formation de la première bulle de vapeur), tandis que la partie droite de de la courbe de saturation est la courbe de rosée  (formation de la première goutte de liquide).


\begin{figure}[H]
\centering
\includegraphics[width=\textwidth]{ph_eau_bis.pdf}
\caption{Diagramme $(p,h)$ pour l'eau (R718). La pression $p$ est en ordonnée ; l'enthalpie massique $h$ en abscisse. Attention à l'échelle logarithmique sur l'axe des ordonnées (pression $p$) !}
\label{fig:frigoristes_eau}
\end{figure}

\parag{Réseaux de courbes} Le diagramme présente plusieurs réseaux de courbes. Les courbes les plus importantes sont les suivantes.
\begin{liste}
\item Les isobares $(p=\cte$, valeur de $p$ en bar) correspondent à des lignes horizontales.
\item Les isenthalpes $(h=\cte$, valeur de $h$ en $\SI{}{kJ.kg^{-1}})$ correspondent à des lignes verticales.
\item  Les \textcolor{JoliRouge}{isothermes}  $(T=\cte$, valeur de $T$ en $\SI{}{\celsius})$.
\begin{remarque}
  Les isothermes s’approchent de droites verticales dans les domaines vapeur et liquide à basse pression, ce qui est cohérent avec la limite du gaz parfait et du liquide incompressible où $h=c_pT$. Elles prennent la forme de paliers (droites horizontales, confondues avec les isobares) dans le domaine de coexistence liquide-vapeur. Les isothermes se courbent fortement voisinage du point critique.
\end{remarque}
\item  Les \textcolor{blue}{isentropes} $(s=\cte$, valeur de $s$ en $\SI{}{kJ.K^{-1}.kg^{-1}})$.
\begin{remarque}
Dans la limite du gaz parfait, les isentropes vérifient $T^\gamma p^{1-\gamma}=\cte$ d'après la loi de Laplace. Or $h=c_p T$ et donc  $h^\gamma p^{1-\gamma}=\cte$ et les isentropes ont pour équation $p \propto h^{\gamma/(\gamma-1)}$.
%Dans les autres domaines du diagramme, les isentropes ont une forme plus compliquée.
\end{remarque}
\end{liste}
On représente aussi parfois d'autres types de courbes.
\begin{liste}
\item Les \textbf{isotitres} ($x_v=\cte$, sans unité) n’existent que dans le domaine de coexistence liquide-vapeur, et indiquent le titre massique en vapeur $x_v = \frac{m_v}{m_\ell + m_v}$. Les isotitres se rejoignent toutes au point critique. On peut déterminer le titre massique en vapeur en un poirte quelconque $M$ sous la courbe de saturation avec la loi des moments.
% On repère également à gauche une partie de la courbe d'ébullition, et à droite une portion de la courbe de rosée (trais noir épais).
\item Les \JVert{isochores} $(v=\cte$, valeur de $v$ en $\SI{}{m^2.kg^{-1}}$).
\end{liste}


\begin{methode}[Utiliser un diagramme $(p,h)$ pour étudier un cycle industriel]
  \item Bien identifier les courbes présentes sur le diagramme $(p,h)$ (saturation, isotherme, etc.).
  \item Placer les points du cycle sur le diagramme $(p,h)$ en s'aidant des données : pression $p$, température $T$, et nature des transformations (isenthalpique, isobare, isentropique, etc.).
  \item Identifier les étapes de contacts avec une source chaude ou froide, et celles où le fluide travaille.
  \item Calculer $\Delta h $ pour ces étapes et l'identifier à $q$ ou $w_u$ en invoquant le
  $1^{\rm er}$ principe industriel.
\end{methode}


\begin{exemple}{Cycle réfrigérant au CO$_2$}
  % Le CO$_2$ est de plus en plus utilisé pour la réfrigération et la climatisation. En effet, on peut l'extraire facilement de l’environnement, et son empreinte environnementale est faible par rapport aux autres fluides frigorigènes.
  On étudie ici le cycle de réfrigération au CO$_2$  suivant : une compression adiabatique réversible dans un compresseur, d'un état 1 ($p_1=\SI{35}{bar}$, température $T_1=\SI{0}{\dC}$, vapeur saturante sèche) à un état $2$ de pression $p_2=\SI{90}{bar}$; un refroidissement isobare dans un condenseur de l'état $2$ à un état $3$ de température $T_3=\SI{40}{\celsius}$;
  une détente isenthalpe dans un détendeur de l'état $3$ à l'état $4$ de pression $p_4=p_1$; un échauffement isobare dans un évaporateur de l'état $4$ jusqu'à retrouver l'état initial $1$.
  \begin{fig}[CO2]
  \includegraphics[width=\textwidth]{R744_CO2_P_h.pdf}
  \capt{Diagramme $(p,h)$ du CO$_2$.}
  \end{fig}
  %%%%%%%%%%%%%
  %\caption{Diagramme $(p,h)$ du CO$_2$.}
  %\label{fig:CO2}
  %\end{figure}
  \begin{quest}
  \item Dessiner le cycle orienté sur le diagramme enthalpique de la \Fig{CO2}. Que vaut la température $T_2$ ? La pression $p_3$ ? Quel est l'état du fluide en $4$ ? Que vaut la température $T_4$ ?

  \item Calculer le transfert thermique massique $q_f$ soutiré à la source froide et le travail utile $w_u$ fourni au fluide par le compresseur.
  % \item Calculer l'efficacité frigorifique du cycle. Comment l'améliorer?
  \end{quest}
\end{exemple}
\examplesol{
\setcounter{quest}{0}
\begin{quest}
\item
Le fluide est diphasé en $4$ et $T_4=T_1=\SI{0}{\dC}$ car une isotherme est confondue avec une isobare  sur le palier de transition. On lit $T_2 = \SI{75}{\dC}$ sur les isothermes (en rouge). La pression en $3$ vaut $p_3 = p_2 = \SI{90}{bar}$ puisque la transformation $2\rightarrow3$ est isobare.
\begin{center}
\includegraphics[width=0.85\textwidth]{phCO2_soluce.pdf}
\end{center}


\item
Il s'agit ici d'un cycle diphasé : le gaz carbonique est partiellement liquéfié lorsqu'il entre en contact avec la source froide, puis il se vaporise entièrement en lui soutirant de l'énergie thermique.
Le contact avec la source froide s'effectue lors de l'échauffement isobare $4\rightarrow 1$, dans l'évaporateur. Le transfert thermique reçu vaut
\eqbox{q_f = \Delta h = h_1-h_4 = 430-325 = \SI{105}{kJ.kg^{-1}}.}
 Le travail utile est fourni au gaz dans le compresseur (phase $1\rightarrow 2$), soit
 \eqbox{w_u = \Delta h =h_2-h_1 = 470-430=\SI{50}{kJ.kg^{-1}}.}

\end{quest}
}

\subsection{Diagramme $(T,s)$}
\label{sec:Ts}

 Un diagramme entropique représente la température $T$ en ordonnée en échelle linéaire et l’entropie $s$ en abscisse en échelle linéaire. Un exemple est le  diagramme de l’eau (réfrigérant R718) de la \Fig{T_s_eau}. Analysons ce diagramme.


\begin{figure}[h!]
\centering
\includegraphics[width=\textwidth]{Ts_eau_bis.pdf}
\caption{Diagramme $(T,s)$ pour l'eau (R718).}
\label{fig:T_s_eau}
\end{figure}


\parag{Courbe de saturation} La courbe en cloche correspond à la courbe de saturation, mais elle n’a pas la même forme que dans le diagramme $(p,h)$ !  Sous la courbe de saturation, le fluide est diphasé : on a coexistence liquide-vapeur $L+V$. De part et d'autre de la courbe de saturation, le fluide est sous une seule phase : liquide $L$ aux basses entropies, et vapeur $V$ aux hautes entropies.
%Le diagramme peut représenter l’eau dans son état liquide $L$ aux basses enthalpies, vapeur $V$ aux hautes enthalpies, ou en coexistence liquide-vapeur $L+V$. La courbe en cloche est la courbe de saturation, qui délimite ces trois régions. Son sommet est le point critique $C$. La partie à gauche de $C$ est la courbe d’ébullition (formation de la première bulle de vapeur) est la partie droite de $C$ est la courbe de rosée  (formation de la première goutte de liquide).

\parag{Réseaux de courbes} Le diagramme présente plusieurs faisceaux de courbes particulières. Les plus importantes sont les suivantes.
\begin{liste}
\item Les isothermes $(T=\cte$, valeur de $T$ en $\SI{}{\dC})$ correspondent à des lignes horizontales.
\item Les isentropes $(s=\cte$, valeur de $s$ en $\SI{}{kJ.K^{-1}.kg^{-1}})$ correspondent à des lignes verticales.
\item  Les \textcolor{JoliRouge}{isobares}  $(p=\cte$, valeur de $p$ en bar).
\begin{remarque}
  Dans le domaine vapeur à basse pression, les isobares ont une allure exponentielle, ce qui est cohérent car dans la limite du gaz parfait à $p=\cte$, $s(T) = s_0 + c_p \ln(T/T_0)$, et donc $T \propto e^{s}$. Elles
  prennent la forme de paliers (segment horizontal, confondu avec les isothermes) dans le domaine de coexistence liquide-vapeur. Les isobares se courbent fortement au voisinage du point critique.
\end{remarque}
 % sont croissantes avec l’entropie massique $s$.
% s’approchent de droites verticales dans le domaine vapeur (cohérent avec la limite du gaz parfait où $h=c_pT$) et dans le domaine liquide (cohérent avec la limite du liquide incompressible où $h=cT$).

\item  Les \textcolor{blue}{isenthalpes} $(h=\cte$, valeur de $h$ en $\SI{}{kJ.K.kg^{-1}})$.
\begin{remarque}
  Les isenthalpes s’approchent  de droites
  horizontales dans le domaine de haute entropie, ce qui est cohérent avec la limite du gaz parfait où $h=c_p T$.
\end{remarque}
%sont indiquées en bleu.de droites verticales dans le domaine liquide (cohérent avec le caractère incompressible du liquide).
\end{liste}

On représente aussi parfois d'autres types de courbes.
\begin{liste}
\item Les \textbf{isotitres} ($x_v=\cte$, sans unité) n’existent que dans le domaine de coexistence liquide-vapeur, et indiquent le titre massique en vapeur $x_v = \frac{m_v}{m_\ell + m_v}$. Les isotitres se rejoignent toutes au point critique. On peut déterminer le titre massique en vapeur en un pointe quelconque $M$ sous la courbe de saturation avec la loi des moments.
% On repère également à gauche une partie de la courbe d'ébullition, et à droite une portion de la courbe de rosée (trais noir épais).
\item Les \JVert{isochores} $(v=\cte$, valeur de $v$ en $\SI{}{m^2.kg^{-1}}$).
\end{liste}

\begin{methode}[Utiliser un diagramme $(T,s)$ pour étudier un cycle industriel]
  \item Bien identifier les courbes présentes sur le diagramme $(T,s)$ (saturation, isenthalpe, etc.).
  \item Placer les points du cycle sur le diagramme $(T,s)$ en s'aidant des données : température $T$, pression $p$, et nature des transformations (isenthalpique, isobare, isentropique, etc.).
  \item Lire graphiquement les valeurs de $h$ en chaque point du cycle en s'aidant des isenthalpes.
  \item Identifier les étapes de contacts avec une source chaude ou froide, et celles où le fluide travaille.
  \item Calculer $\Delta h$ pour ces étapes et l'identifier à $q$ ou $w_u$ en invoquant le
  $1^{\rm er}$ principe industriel.
\end{methode}

\end{document}
